\section{East African prior-predictive calibration}\label{east-african-prior-predictive-calibration}

\subsection{1. Calibration status}\label{calibration-status}

The numerical results form a \textbf{prior-predictive calibration}.

It reports the outcome distribution implied by the current assumptions and
uncertainty distributions. The regional tree, log, and matching operators still
require identification and posterior-predictive validation.

\subsection{2. Evidence used as scale anchors}\label{evidence-used-as-scale-anchors}

The working paper reports approximately:

\begin{itemize}
\tightlist
\item
  338,555 planted hectares in Uganda;
\item
  588,000 planted hectares in Tanzania;
\item
  about 0.93 million hectares combined;
\item
  37 Chinese-owned processors in the internal intelligence base;
\item
  about 5 million m3/year installed capacity;
\item
  about 3 million m3/year current use;
\item
  typical utilisation around 50-60\%.
\end{itemize}

These figures are 2025--2026 evidence inputs cited by the paper. They retain the
status and reference date of their underlying sources.

The paper also uses a Uganda CFR register with about 94,820 recorded planted
hectares across 3,294 positive planted portfolios. The reported planted-area
Gini is approximately 0.803. A single average estate size therefore gives a poor
mapping from hectares to account count.

\subsection{3. Missing structural state}\label{missing-structural-state}

The current calibration lacks a region-wide joint tree-level size-quality
distribution. Regional data therefore leave the full mapping

\[
\mu_t\to\widetilde\Lambda_t\to\Gamma_t^*
\]

unidentified. The reduced model uses priors for commercially available flow,
processor coverage, adoption, monetisation, and correlated market state.

This should be read as a placeholder architecture with explicit destinations
for future evidence.

\subsection{4. Prior-predictive scale results}\label{prior-predictive-scale-results}

The paper reports 50,000 joint prior-predictive draws. The draw count controls
Monte Carlo error conditional on the model. Prior validation and structural
discrepancy require observed outcomes.

Key reported values are:

{\def\LTcaptype{none} % do not increment counter
\begin{longtable}[]{@{}lrrr@{}}
\toprule\noalign{}
Metric & 100k ha & 500k ha & 1m ha \\
\midrule\noalign{}
\endhead
\bottomrule\noalign{}
\endlastfoot
Revenue P20, USDk/yr & 49 & 326 & 650 \\
Revenue P50, USDk/yr & 78 & 451 & 850 \\
Revenue P80, USDk/yr & 115 & 613 & 1,090 \\
Base OPEX P50, USDk/yr & 219 & 424 & 669 \\
Robust OPEX P50, USDk/yr & 350 & 617 & 875 \\
Joint operating surplus P20, USDk/yr & -197 & -125 & -66 \\
Joint operating surplus P50, USDk/yr & -145 & 21 & 171 \\
Joint operating surplus P80, USDk/yr & -97 & 191 & 436 \\
P(operating surplus \textgreater{} 0) & 0.01 & 0.55 & 0.72 \\
Paying-ha penetration & 0.25 & 0.26 & 0.262 \\
Asset fee, USD/paying ha & 2.16 & 2.17 & 2.18 \\
Supply fee, USD/visible t & 0.253 & 0.259 & 0.260 \\
\end{longtable}
}

These outputs are conditional on the paper's mapping from coverage to commercial
activity and carry no causal interpretation for hectares.

\subsection{5. Cost identification example}\label{cost-identification-example}

The three median base-cost anchors are approximately

\[
(100k,219k),\quad(500k,424k),\quad(1m,669k).
\]

A naive pure power model

\[
C(H)=aH^\beta
\]

fits these points with an apparent exponent near \texttt{0.47}.

But the paper gives an exact three-point fit of the form

\[
C(H)=162.8+56.2\left(\frac{H}{100000}\right)^{0.955}
\]

in USD thousands per year.

The variable exponent is then close to one. The low apparent total-cost
elasticity comes largely from spreading the fixed USD 162.8k component.

This is an important model-audit example: a scaling curve can have very
different structural interpretations that agree at the observed anchors.

\subsection{6. Break-even thresholds}\label{break-even-thresholds}

Holding the other revenue side fixed at its current median, the working paper
reports one-lever deterministic break-even thresholds.

{\def\LTcaptype{none} % do not increment counter
\begin{longtable}[]{@{}
  >{\raggedright\arraybackslash}p{(\linewidth - 8\tabcolsep) * \real{0.1667}}
  >{\raggedright\arraybackslash}p{(\linewidth - 8\tabcolsep) * \real{0.1667}}
  >{\raggedleft\arraybackslash}p{(\linewidth - 8\tabcolsep) * \real{0.2222}}
  >{\raggedleft\arraybackslash}p{(\linewidth - 8\tabcolsep) * \real{0.2222}}
  >{\raggedleft\arraybackslash}p{(\linewidth - 8\tabcolsep) * \real{0.2222}}@{}}
\toprule\noalign{}
\begin{minipage}[b]{\linewidth}\raggedright
Scale
\end{minipage} & \begin{minipage}[b]{\linewidth}\raggedright
Cost mode
\end{minipage} & \begin{minipage}[b]{\linewidth}\raggedleft
Asset monetisation (\% managed spend)
\end{minipage} & \begin{minipage}[b]{\linewidth}\raggedleft
Supply fee (USD/t)
\end{minipage} & \begin{minipage}[b]{\linewidth}\raggedleft
Paying-ha penetration
\end{minipage} \\
\midrule\noalign{}
\endhead
\bottomrule\noalign{}
\endlastfoot
100k & Base & 2.23 & 2.04 & 90.3\% \\
100k & Robust & 3.73 & 3.70 & 150.9\% \\
500k & Base & 0.56 & 0.21 & 23.5\% \\
500k & Robust & 0.98 & 0.54 & 41.3\% \\
1m & Base & 0.42 & 0.08 & 17.9\% \\
1m & Robust & 0.65 & 0.29 & 27.3\% \\
\end{longtable}
}

A penetration threshold above 100\% shows that the selected lever has no feasible
break-even point under those assumptions.

\subsection{7. Robustness multiplier}\label{robustness-multiplier}

For an 80\% target probability of nonnegative recurring operating surplus, the
paper reports revenue multipliers of approximately

\[
M_{100k}^{0.8}=4.68,\qquad
M_{500k}^{0.8}=1.37,\qquad
M_{1m}^{0.8}=1.10.
\]

This measures distance from the target under the current joint distributions.
It supplies no causal effect of scale.

\subsection{8. Adoption priors}\label{adoption-priors}

The appendix gives size-dependent asset-adoption priors at central price/value
conditions:

{\scriptsize\def\LTcaptype{none} % do not increment counter
\begin{longtable}[]{@{}
  >{\raggedright\arraybackslash}p{(\linewidth - 14\tabcolsep) * \real{0.0968}}
  >{\raggedleft\arraybackslash}p{(\linewidth - 14\tabcolsep) * \real{0.1290}}
  >{\raggedleft\arraybackslash}p{(\linewidth - 14\tabcolsep) * \real{0.1290}}
  >{\raggedleft\arraybackslash}p{(\linewidth - 14\tabcolsep) * \real{0.1290}}
  >{\raggedleft\arraybackslash}p{(\linewidth - 14\tabcolsep) * \real{0.1290}}
  >{\raggedleft\arraybackslash}p{(\linewidth - 14\tabcolsep) * \real{0.1290}}
  >{\raggedleft\arraybackslash}p{(\linewidth - 14\tabcolsep) * \real{0.1290}}
  >{\raggedleft\arraybackslash}p{(\linewidth - 14\tabcolsep) * \real{0.1290}}@{}}
\toprule\noalign{}
\begin{minipage}[b]{\linewidth}\raggedright
Portfolio size
\end{minipage} & \begin{minipage}[b]{\linewidth}\raggedleft
1-10
\end{minipage} & \begin{minipage}[b]{\linewidth}\raggedleft
10-50
\end{minipage} & \begin{minipage}[b]{\linewidth}\raggedleft
50-150
\end{minipage} & \begin{minipage}[b]{\linewidth}\raggedleft
150-500
\end{minipage} & \begin{minipage}[b]{\linewidth}\raggedleft
500-1,500
\end{minipage} & \begin{minipage}[b]{\linewidth}\raggedleft
1,500-5,000
\end{minipage} & \begin{minipage}[b]{\linewidth}\raggedleft
5,000+
\end{minipage} \\
\midrule\noalign{}
\endhead
\bottomrule\noalign{}
\endlastfoot
Base adoption & 0.1\% & 0.5\% & 5\% & 15\% & 35\% & 60\% & 80\% \\
\end{longtable}
}

These values are priors. Observed conversion and renewal data will replace them.

\subsection{9. What new evidence would structurally improve the calibration?}\label{what-new-evidence-would-structurally-improve-the-calibration}

The model identifies specific targets:

\begin{itemize}
\tightlist
\item
  repeated tree/plot measurements for growth, mortality, and quality transition;
\item
  harvest-linked tree-to-log records for recovery and bucking;
\item
  operator logs for productivity and delivered cost;
\item
  processor intake for grade acceptance and density conversion;
\item
  repeated bids/alternative buyers for outside options and bargaining;
\item
  observed platform conversion and renewal for willingness to pay.
\end{itemize}

Each dataset has a specified target operator and likelihood.

\subsection{10. Interpretation discipline}\label{interpretation-discipline}

The empirical programme follows this sequence:

\begin{Shaded}
\begin{Highlighting}[]
\NormalTok{prior{-}predictive scale experiment}
\NormalTok{    {-}\textgreater{} targeted data collection}
\NormalTok{    {-}\textgreater{} operator identification}
\NormalTok{    {-}\textgreater{} posterior predictive validation}
\NormalTok{    {-}\textgreater{} decision experiments}
\NormalTok{    {-}\textgreater{} stronger operational claims}
\end{Highlighting}
\end{Shaded}

The calibration is therefore a research instrument as much as a business-case
output.
